\documentclass[10pt,a4paper]{amsart}
\usepackage{amsmath,amssymb,amsthm,graphicx}
% Retained minimal preamble for the Stacks Project chapter more-morphisms.tex (More on Morphisms).
% Verbatim excerpt of the chapter preamble: Xy-pic setup (native preamble lines 14-18) and the native
% theorem declarations with the mathematical macros (native preamble lines 154-217). Copied mechanically;
% nothing added, removed or redefined. The wrapper class and the amsmath/amssymb/amsthm/graphicx base
% packages are supplied by the assembler, so the class/externaldocument part of the native preamble is omitted.
% For commutative diagrams we use Xy-pic
\usepackage[all]{xy}

% We use 2cell for 2-commutative diagrams.
\xyoption{2cell}
% Theorem environments.
%
\theoremstyle{plain}
\newtheorem{theorem}[subsection]{Theorem}
\newtheorem{proposition}[subsection]{Proposition}
\newtheorem{lemma}[subsection]{Lemma}

\theoremstyle{definition}
\newtheorem{definition}[subsection]{Definition}
\newtheorem{example}[subsection]{Example}
\newtheorem{exercise}[subsection]{Exercise}
\newtheorem{situation}[subsection]{Situation}

\theoremstyle{remark}
\newtheorem{remark}[subsection]{Remark}
\newtheorem{remarks}[subsection]{Remarks}

\numberwithin{equation}{subsection}

% Macros
%
\def\lim{\mathop{\mathrm{lim}}\nolimits}
\def\colim{\mathop{\mathrm{colim}}\nolimits}
\def\Spec{\mathop{\mathrm{Spec}}}
\def\Hom{\mathop{\mathrm{Hom}}\nolimits}
\def\Ext{\mathop{\mathrm{Ext}}\nolimits}
\def\SheafHom{\mathop{\mathcal{H}\!\mathit{om}}\nolimits}
\def\SheafExt{\mathop{\mathcal{E}\!\mathit{xt}}\nolimits}
\def\Sch{\mathit{Sch}}
\def\Mor{\mathop{\mathrm{Mor}}\nolimits}
\def\Ob{\mathop{\mathrm{Ob}}\nolimits}
\def\Sh{\mathop{\mathit{Sh}}\nolimits}
\def\NL{\mathop{N\!L}\nolimits}
\def\CH{\mathop{\mathrm{CH}}\nolimits}
\def\proetale{{pro\text{-}\acute{e}tale}}
\def\etale{{\acute{e}tale}}
\def\QCoh{\mathit{QCoh}}
\def\Ker{\mathop{\mathrm{Ker}}}
\def\Im{\mathop{\mathrm{Im}}}
\def\Coker{\mathop{\mathrm{Coker}}}
\def\Coim{\mathop{\mathrm{Coim}}}

% Boxtimes
%
\DeclareMathSymbol{\boxtimes}{\mathbin}{AMSa}{"02}

%
% Macros for moduli stacks/spaces
%
\def\QCohstack{\mathcal{QC}\!\mathit{oh}}
\def\Cohstack{\mathcal{C}\!\mathit{oh}}
\def\Spacesstack{\mathcal{S}\!\mathit{paces}}
\def\Quotfunctor{\mathrm{Quot}}
\def\Hilbfunctor{\mathrm{Hilb}}
\def\Curvesstack{\mathcal{C}\!\mathit{urves}}
\def\Polarizedstack{\mathcal{P}\!\mathit{olarized}}
\def\Complexesstack{\mathcal{C}\!\mathit{omplexes}}
% \Pic is the operator that assigns to X its picard group, usage \Pic(X)
% \Picardstack_{X/B} denotes the Picard stack of X over B
% \Picardfunctor_{X/B} denotes the Picard functor of X over B
\def\Pic{\mathop{\mathrm{Pic}}\nolimits}
\def\Picardstack{\mathcal{P}\!\mathit{ic}}
\def\Picardfunctor{\mathrm{Pic}}
\def\Deformationcategory{\mathcal{D}\!\mathit{ef}}
\newenvironment{reference}{}{}
\newenvironment{slogan}{}{}
\newenvironment{history}{}{}
\usepackage{cleveref}
\crefname{originallemma}{original lemma}{original lemmas}
\Crefname{originallemma}{Original lemma}{Original lemmas}
\makeatletter
\newlabel{lemma-etale-makes-quasi-finite-finite-at-point}{{02LK}{source}{}{sourcepointer-1}{}}
\newlabel{lemma-etale-makes-quasi-finite-finite-at-point@cref}{{[originallemma][0][]02LK}{[1][0][]source}{}{}{}}
\makeatother
\title{Making several isolated fibre points finite over one elementary etale neighbourhood}
\author{The Stacks Project Authors}
\date{stacks-project:more-morphisms.tex:local-snapshot-20261011; GNU FDL 1.2 or later}
\begin{document}
\maketitle
\noindent\textit{Editorial note.} Statement span is lines 11973-11995 of sources/stacks-more-morphisms/pinned-original-source.tex; proof span is lines 11997-12012, the complete author proof (14 body lines).\par\medskip
\noindent\textit{Editorial note.} Navigation stubs: 1 cited-but-not-included labels are supplied as non-load-bearing pointers with their pinned tags where the snapshot tag index has them.\par\medskip
\noindent\textit{Editorial note.} extra\_packages is empty on purpose: the retained preamble already loads xy with options and a second declaration clashes in the assembler wrapper.\par\medskip
\noindent\textit{Editorial note.} Text is retained verbatim, including the author's own typos and placeholders.\par\medskip
\noindent\textit{Allowed standard prerequisites (editorial).}\begin{itemize}
\item Explicit prerequisite boundary. Elementary etale neighbourhoods, isolated points of a fibre, and finiteness of a morphism over an etale base change; the assembly is retained.
\end{itemize}\medskip
\ifdefined\hypertarget\hypertarget{sourcepointer-1}{}\fi
\noindent\textit{Source reference.} Original lemma 02LK: sources/stacks-more-morphisms/pinned-original-source.tex, source lines 11924--11934. Cited by the retained statement or proof but not part of the collected spans; supplied only as a navigation pointer. The argument does not depend on its content.\par\medskip
\begin{lemma}
\label{lemma-etale-makes-quasi-finite-finite-multiple-points}
Let $f : X \to S$ be a morphism of schemes.
Let $x_1, \ldots, x_n \in X$ be points having the same image $s$ in $S$.
Assume that
\begin{enumerate}
\item $f$ is locally of finite type, and
\item $x_i \in X_s$ is isolated for $i = 1, \ldots, n$.
\end{enumerate}
Then there exist
\begin{enumerate}
\item[(a)] an elementary \'etale neighbourhood $(U, u) \to (S, s)$,
\item[(b)] for each $i$ an open subscheme $V_i \subset X_U$,
\end{enumerate}
such that for each $i$ we have
\begin{enumerate}
\item[(\romannumeral1)] $V_i \to U$ is a finite morphism,
\item[(\romannumeral2)] there is a unique point $v_i$ of $V_i$
mapping to $u$ in $U$, and
\item[(\romannumeral3)] the point $v_i$ maps to $x_i$ in $X$ and
$\kappa(x_i) = \kappa(v_i)$.
\end{enumerate}
\end{lemma}

\begin{proof}
We will use induction on $n$.
Namely, suppose $(U, u) \to (S, s)$ and $V_i \subset X_U$,
$i = 1, \ldots, n - 1$ work for $x_1, \ldots, x_{n - 1}$. Since
$\kappa(s) = \kappa(u)$ the fibre $(X_U)_u = X_s$. Hence there
exists a unique point $x'_n \in X_u \subset X_U$ corresponding to
$x_n \in X_s$. Also $x'_n$ is isolated in $X_u$. Hence by
Lemma \ref{lemma-etale-makes-quasi-finite-finite-at-point} there
exists an elementary \'etale neighbourhood $(U', u') \to (U, u)$
and an open $V_n \subset X_{U'}$ which works for $x'_n$ and hence
for $x_n$.
By the final assertion of
Lemma \ref{lemma-etale-makes-quasi-finite-finite-at-point}
the open subschemes $V'_i = U'\times_U V_i$ for $i = 1, \ldots, n - 1$ still
work with respect to $x_1, \ldots, x_{n - 1}$. Hence we win.
\end{proof}

\end{document}
